% ============================================================
% PyQCD: tag dev25 至 2026-10-25 HEAD 的胶子 TMD 物理结果审计
% XeLaTeX, two passes. Source and PDF share the same stem.
% ============================================================
\documentclass[UTF8]{ctexart}
\usepackage[margin=2.5cm]{geometry}
\usepackage{graphicx}
\usepackage{amsmath}
\usepackage{microtype}
\usepackage{listings}
\usepackage{xcolor}
\usepackage{booktabs}
\usepackage{longtable}
\usepackage{xltabular}
\usepackage{tabularx}
\usepackage{hyperref}
\usepackage{fancyhdr}
\usepackage[most]{tcolorbox}
\usepackage{titlesec}

\definecolor{accent}{HTML}{1F4E79}
\definecolor{evidence}{HTML}{1E8449}
\definecolor{reference}{HTML}{8C6D1F}
\definecolor{conclusion}{HTML}{8B1A1A}
\definecolor{codebg}{HTML}{F4F6F7}
\definecolor{struct}{HTML}{3B3B98}
\definecolor{relation}{HTML}{0E7C7B}
\definecolor{idea}{HTML}{B03A2E}
\definecolor{warn}{HTML}{D35400}
\definecolor{hl}{HTML}{FDF6E3}

\setlength{\emergencystretch}{3em}
\setlength{\headheight}{13pt}
\hypersetup{
  colorlinks=true,
  linkcolor=accent,
  urlcolor=relation,
  pdftitle={PyQCD dev25 至 2026-10-25 HEAD 胶子 TMD 物理结果审计},
  pdfauthor={PyQCD physics and code audit}
}

\titleformat{\section}{\Large\bfseries\color{accent}}{\thesection}{1em}{}
\titleformat{\subsection}{\large\bfseries\color{struct}}{\thesubsection}{1em}{}
\titleformat{\subsubsection}{\normalsize\bfseries\color{struct!70!black}}{\thesubsubsection}{1em}{}

\newtcolorbox{conclbox}{breakable,colback=conclusion!6,colframe=conclusion,title=结论}
\newtcolorbox{refbox}{breakable,colback=reference!6,colframe=reference,title=参考背景}
\newtcolorbox{ideabox}{breakable,colback=idea!6,colframe=idea,title=物理与工程判断}
\newtcolorbox{warnbox}{breakable,colback=warn!6,colframe=warn,title=局限与风险}
\newtcolorbox{keybox}{breakable,colback=hl,colframe=accent,title=要点}

\lstset{
  basicstyle=\ttfamily\footnotesize,
  backgroundcolor=\color{codebg},
  frame=single,
  language=Python,
  firstnumber=auto,
  keywordstyle=\color{accent}\bfseries,
  commentstyle=\color{evidence!70!black},
  stringstyle=\color{struct},
  breaklines=true,
  breakatwhitespace=false,
  postbreak=\mbox{\textcolor{accent}{$\hookrightarrow$}\space},
  keepspaces=true,
  columns=flexible,
  numbers=left,
  numberstyle=\tiny\color{struct!60!black},
}

\newcommand{\code}[1]{\texttt{\detokenize{#1}}}
\newcommand{\pathref}[1]{\texttt{\nolinkurl{#1}}}
\newcommand{\MSbar}{\overline{\mathrm{MS}}}
\newcommand{\GF}{\mathrm{GF}}
\newcommand{\Tr}{\operatorname{Tr}}
\newcommand{\Real}{\operatorname{Re}}
\newcommand{\Imag}{\operatorname{Im}}
\newcommand{\dd}{\mathrm{d}}
\newcommand{\bT}{\boldsymbol b_T}
\newcommand{\kT}{\boldsymbol k_T}
\newcommand{\Pz}{P_z}
\newcommand{\ts}{t_{\mathrm{sep}}}
\newcommand{\statusdiag}{\textcolor{warn}{\textbf{diagnostic\_only}}}
\newcommand{\confirmed}{\textcolor{struct}{\textbf{确证}}}
\newcommand{\inferred}{\textcolor{relation}{\textbf{推断}}}
\newcommand{\unverified}{\textcolor{warn}{\textbf{未验证}}}

\pagestyle{fancy}
\fancyhf{}
\fancyhead[L]{PyQCD: dev25 至 2026-10-25 HEAD 物理结果审计}
\fancyhead[R]{\thepage}
\renewcommand{\headrulewidth}{0.4pt}

\title{PyQCD 从 \texttt{dev25} 到当前 HEAD 的\\胶子 TMD 物理结果、公式与图表审计}
\author{PyQCD 物理链审计}
\date{2026 年 10 月 25 日}

\begin{document}
\maketitle

\begin{abstract}
本报告审计 PyQCD 从标签 \texttt{dev25}（annotated tag 对象
\texttt{0840298}，peeled commit \texttt{2d21a74}）之后至 2026-10-25 当前 HEAD
\texttt{0f5ac61} 的物理与算法增量。该区间共 14 个提交，实体成果集中在
\texttt{dev27}（提交 \texttt{19ab351}）：有限横距 Wilson staple、
full-channel 胶子 OPE、source-resolved \texttt{nopol/pol35} 核子两点函数、
一维与二维 ratio-v3、delete-one jackknife 协方差、二维 staircase 顺序审计，
以及一次 11 组态真实数据诊断。软件侧同时修复了 GPU Wilson-flow 六-staple
峰值显存缺陷，当前全量回归记录为 112 项通过。

物理上，本区间已经建立了可复核的
\[
\text{流化规范场}\longrightarrow\text{有限 }b_\perp\text{ 算符}
\longrightarrow C_2/C_3\longrightarrow
\text{ratio/JK 协方差}
\]
诊断链，并在 4150--8050 的 11 个组态上运行。它尚未建立核子胶子
TMD-PDF 的物理闭环：\texttt{t\_sep} 漂移显著超过 JK 误差；所有非零
\(b_T/a=1,\ldots,4\) 均落在 \(r_f/a=\sqrt{24}\approx4.899\) 的流半径内；
\(\Pz r_f/(\hbar c)\) 对 P200/P400 分别约为 2.57/5.13；并且没有同几何
soft/rapidity、完整胶子-夸克 matching、多流时或多格距连续极限。
因此所有 all11 结果必须保持 \statusdiag 与
\code{physical_acceptance=false}。

本报告同时记录一个必须显式处理的估计量分歧：ratio-v3 使用完整复数
\texttt{nopol} \(C_2\) 作共同分母，而归档 donghx ratio 与仓内主
\texttt{\_disconnected} 路径使用 \(\Real C_{2,\mathrm{nopol}}\)。在
\(\ts/a=8\) 的 source 平均中，P400 的
\(|\Imag C_2/\Real C_2|=15.60\%\)（相位 \(8.87^\circ\)）；这是\textbf{单点}
相位敏感度，不是跨 \(\ts\) 的稳定相位，也不是 ratio 的统计误差。
同一受控复输入的实测 denominator 差为
\(\max|\Delta R|=6.77\times10^{-3}\)（归一化最大差 7.86\%，
相对 \(L_2\) 7.28\%）。在真实 11 组态上，该约定差在 P400、\(\ts/a=10\)
达到 \(\max|\Delta R|/\sigma_{\rm JK}=5.7\)（unpolarized）与 \(3.1\)（helicity），
而 P000/P200 均小于 \(0.2\)。这不是胶子-夸克算符混合，而是估计量/相位约定差异；
在统一之前，不能把当前数值与 donghx 的绝对 ratio 直接等同。
\end{abstract}

\tableofcontents

\section{任务、版本边界与结论先行}

\subsection{任务定义与证据范围}

任务定义为：梳理 \code{tag dev25} 到当前 HEAD 的全部物理结果，提取涉及的物理公式和
相关图表，结合用户指定的参考文件
\pathref{/root/MyQCD/docs/pyqcd_physics/report_gluon_tmd_lattice_qcd_review_20261023.tex}
进行大量物理分析，并在本仓库 \code{docs/} 中生成可编译的 LaTeX 与 PDF。

\begin{table}[htbp]
\centering
\small
\caption{版本边界（来源：\code{git show-ref}、\code{git log dev25..HEAD} 实测）}
\begin{tabularx}{\textwidth}{@{}l l X@{}}
\toprule
对象 & 提交/标签 & 作用 \\
\midrule
审计起点 & \code{0840298}（tag）/ \code{2d21a74}（commit） & \code{dev25} annotated tag 的 2026-10-23 发布快照 \\
中继标签 & \code{4c4b662} / \code{dev26} & 参考脚本、PyQCU 子模块与归档整理 \\
物理实体标签 & \code{19ab351} / \code{dev27} & 11 组态极化 TMD 诊断与二维批量链 \\
当前 HEAD & \code{0f5ac61} & \code{dev27} 之后 4 个归档/下载状态提交 \\
提交数量 & 14 & \code{dev25..HEAD} \\
工作树 & 派生输出未跟踪 & 报告、图、构建与审计文件均为 HEAD 之后新增；历史“干净”状态无独立快照，标为未验证 \\
\bottomrule
\end{tabularx}
\end{table}

\subsection{证据分级}

\begin{table}[htbp]
\centering
\small
\caption{本报告采用的证据等级（来源：代码、日志、运行产物与参考理论交叉核对）}
\begin{tabularx}{\textwidth}{@{}p{0.22\textwidth}X@{}}
\toprule
状态 & 含义与本次实例 \\
\midrule
\confirmed{} 算符/代码 &
函数可定位，几何、shape、dtype、路径和代数恒等式由源码及测试确定 \\
\confirmed{} 真实材料 &
11 组态 C2/OPE/ratio 文件存在，有限性、索引哈希和 JK 样本可直接读取 \\
\confirmed{} 软件回归 &
当前 HEAD 的测试和编译命令实际返回成功；只说明软件契约，不说明物理闭环 \\
\inferred{} 物理解读 &
把有限 \(b_T\)、\(\ts\)、\(\Pz\) 诊断与 TMD 结构联系起来的解释，必须声明尺度和近似 \\
\unverified{} 物理闭环 &
soft/rapidity、完整 matching、混合矩阵、多尺度极限或稳定基态平台尚无证据 \\
\bottomrule
\end{tabularx}
\end{table}

\subsection{结论先行}

\begin{conclbox}
\begin{enumerate}
  \item \textbf{本区间确实产生了新的真实数据结果。} 11 个组态、3 个动量、
        1D \(b=0,\ldots,4\)、2D \((b_x,b_y)=(0,\ldots,4)\times(0,1)\)、
        6 个 \(\ts/a\) 构成可追溯的 Euclidean 胶子矩阵元与 ratio 诊断集
        \cite{local-log-all11}.
  \item \textbf{算符与合同层已经达到可继续开发的稳定程度。} 有限 \(b\)
        geometry、\(b=0\) 退化、纯规范零响应、双极化 C2、full-channel
        OPE、cache 身份和二维 staircase 顺序均有独立契约或真实数据对照。
  \item \textbf{物理基态矩阵元尚未提取。} 1D unpolarized 的最大
        \(\ts\) 漂移/误差为 7.86、10.17、12.95（P000/P200/P400）；
        2D 为 19.84、8.53、63.57。该指标是
        \(\max_{z,b}[\max_t R-\min_t R]/\overline{\sigma_t}\)，只证明
        \textbf{至少一个} \((z,b)\) 点失败。按朴素判据 \(\le1\)，局部通过率为
        1D \(22/65\)、\(26/65\)、\(13/65\)（unpolarized）与
        \(43/65\)、\(20/65\)、\(16/65\)（helicity）；2D 为
        \(43/130\)、\(45/130\)、\(24/130\) 与 \(82/130\)、\(41/130\)、\(28/130\)。
        因此正确表述是：\textbf{没有任何动量在整个 \((z,b)\) 面上通过均匀平台门}；
        局部通过点仍需相关协方差和预先声明的选点规则检验，当前未做。
  \item \textbf{流窗口不成立。} \(a\simeq0.1053\,\mathrm{fm}\)，
        \(r_f=4.899a\simeq0.516\,\mathrm{fm}\)，而全部非零 1D 位移
        \(b\le4a=0.421\,\mathrm{fm}\)，二维最大模长
        \(|\bT|_{\max}=\sqrt{17}a=0.434\,\mathrm{fm}\)；\(z/a=1,\ldots,4\)
        同样满足 \(z<r_f\)。流化尺度大于待测横向与纵向距离。
  \item \textbf{大动量条件也不成立。} P200/P400 对应约
        \(0.981/1.962\,\mathrm{GeV}\)，而
        \(\Pz r_f/(\hbar c)\simeq2.57/5.13\)（格点乘积
        \((\Pz a)(r_f/a)\) 同值），不满足 \(\Pz r_f\ll1\)。
  \item \textbf{复/实 \texttt{nopol} 分母尚未统一。} ratio-v3 使用复数除法；
        归档 donghx 与仓内主 \texttt{\_disconnected} 使用实数分母。P400 在
        \(\ts/a=8\) 的单点相位为 \(8.87^\circ\)（\(|\Imag/\Real|=15.60\%\)），
        受控复输入的 denominator 差为 \(\max|\Delta R|=6.77\times10^{-3}\)；
        两者都不是 ratio 的统计误差，也不能外推为稳定相位。
  \item \textbf{二维路径顺序是算符身份的一部分。} \(b_y/a=0\) 时
        \((1,0)\) 与 \((0,1)\) 逐元素一致；\(b_y/a=1\) 的切片相对 \(L_2\)
        对 4150/6250 分别为 1.954\% 与 1.793\%（完整二维数组为
        1.156\%/1.076\%）。这是 contour/Wilson-loop 结构，不是浮点舍入。
  \item \textbf{因此当前应称为：} 流化、source-resolved、有限横距胶子
        Euclidean 矩阵元与 disconnected ratio 诊断；不应称为光锥胶子
        TMD-PDF、\(\Delta G(x,\bT)\) 或已经完成梯度流重整化的物理结果。
\end{enumerate}
\end{conclbox}

\section{主体结构、数据关系与项目思路}

\subsection{提交与阶段结构}

\begin{xltabular}{\textwidth}{@{}l l X@{}}
\caption{dev25 之后 14 个提交的阶段归并（数据来源：Git 实测）}\\
\toprule
提交 & 阶段 & 物理/工程内容 \\
\midrule
\endfirsthead
\toprule
提交 & 阶段 & 物理/工程内容 \\
\midrule
\endhead
\code{cdb81c7} & 参考整理 & huangcl P0 ratio 与 sush 2pt 动态缩并参考代码调整；不是 PyQCD 生产结果 \\
\code{4c4b662} & 子模块 & 新增 PyQCU 子模块，无物理数值结论 \\
\code{158d887} & dev26 日志 & 归档 dev26 核查；记录后续空白清理事项 \\
\code{395dba0} & 交叉验证 & 新增 donghx 4D 小格点对照驱动，区分直线准 PDF 与有限 \(b\) TMD 扩展 \\
\code{36de54c} & 核心落库 & 有限 \(b\) geometry、full-channel OPE、双极化 C2、1D/2D cache、ratio-v2/v3、真实输入审计与日志 \\
\code{f432c36} & 参考同步 & donghx/huangcl/sush 参考资料批量同步，无新生产运行 \\
\code{1354d31} & 数据盘点 & 682.023 TB 可读关键数据与 101 项传输清单，工程信息而非物理结果 \\
\code{06826a3} & GPU 修复 & 按需 roll 消除六-staple 和 RK3 同时驻留导致的峰值显存缺陷 \\
\code{7d4e74d} & JK 持久化 & 保存完整 delete-one JK 组态轴，并收紧 complex64 orientation 容差 \\
\code{19ab351} & dev27 & 11 组态真实 C2、1D/2D OPE、audit、ratio、covariance 与 staircase 诊断 \\
\code{815e964} & 会话转储 & 临时加入 89,302 行会话 Markdown，无算法内容 \\
\code{c89ee42} & 转储删除 & 删除上一提交同一文件，净效果为零 \\
\code{ac5a9cf} & 下载标记 & 加入 8 个传输完成标记，无物理影响 \\
\code{0f5ac61} & 忽略规则 & 删除完成标记并忽略下载状态目录 \\
\bottomrule
\end{xltabular}

\subsection{模块与物理对象}

\begin{xltabular}{\textwidth}{@{}p{0.30\textwidth}p{0.30\textwidth}X@{}}
\caption{代码对象与物理对象映射（来源：当前 HEAD 源码）}\\
\toprule
代码入口 & 物理对象 & 本区间作用与边界 \\
\midrule
\endfirsthead
\toprule
代码入口 & 物理对象 & 本区间作用与边界 \\
\midrule
\endhead
\pathref{pyqcd/renorm/_gradient_flow.py} & Wilson flow \(V_\tau\) & RK3 积分；GPU 峰值内存修复，不改变数值约定 \\
\pathref{pyqcd/renorm/_tmd.py} & Clover \(F_{\mu\nu}\)、staple、胶子双线性 & 新增 1D/2D finite-\(b\) full-channel OPE，输出复数逐时间片可观测量 \\
\pathref{pyqcd/pipeline/_steps.py} & 核子 C2 与投影 & 一次 \(Cg5g4\) 合同生成 \texttt{nopol/pol35} source-resolved C2 \\
\pathref{pyqcd/pipeline/_tmd9.py} & canonical cache/loader & 绑定极化、源轴、\(z,b\)、staircase、方向、精度和 SHA-256 \\
\pathref{pyqcd/analysis/_gradient_flow_ratio.py} & disconnected \(C_3/C_2\) & 显式 \(z,b_x,b_y\) 轴、复数估计、delete-one JK、最终 helicity 相位 \\
\pathref{pyqcd/testing/tmd/audit_finite_b_inputs.py} & fail-closed 输入门 & 组态数、双极化 C2、OPE shape/labels/order 与共享索引审计 \\
\pathref{pyqcd/testing/tmd/run_real_polarized_chain.py} & 批量恢复 & 1D/2D、GPU/CPU、cache hit、逐组态 checkpoint 与失败恢复 \\
\bottomrule
\end{xltabular}

\subsection{数据流与身份}

\begin{table}[htbp]
\centering
\small
\caption{从组态到 ratio 诊断的可执行流程（数据来源：源码与真实 run 产物）}
\begin{tabularx}{\textwidth}{@{}l X@{}}
\toprule
阶段 & 关键输入、输出与守恒量 \\
\midrule
输入 & 真实 gauge、perambulator、VVV；\(24^3\times72\)，\(a\simeq0.1053\,\mathrm{fm}\) \\
流化 & \(\tau=3\)、\(\epsilon=0.05\)；输出 flowed links \(V_\tau\) \\
OPE & \(z/a=0,\ldots,12\)；1D \(b/a=0,\ldots,4\)；2D \(b_x/a=0,\ldots,4,\ b_y/a=0,1\) \\
外态 & \texttt{nopol/pol35} C2，形状 \((72,72)\)，时间轴 source/relative-$dt$ \\
审计 & 共享组态顺序、输入索引 SHA、极化成对覆盖、schema、shape、dtype、finite \\
ratio & \(C_3/C_2\)、维度 \((z,b_x,b_y)\times\ts\times\) insertion \(\times P_z\) \\
统计 & delete-one JK 共 11 个样本；协方差最大正秩 \(N_{\rm conf}-1=10\) \\
物理门 & 当前止于 plateau/流窗口/soft/rapidity/matching/continuum 之前 \\
\bottomrule
\end{tabularx}
\end{table}

\subsection{设计意图}

\begin{ideabox}
该区间的工程设计不是“尽快得到一条 PDF 曲线”，而是先建立可追踪的中间物理对象：
同一组态上的 OPE 与 C2、来源明确的极化和源轴、显式离散路径顺序、不可混淆的
cache 身份，以及 fail-closed 的 audit。这个顺序是正确的，因为胶子 TMD 的
Wilson-line 几何、算符混合和软/快度结构都不能由“文件存在”或“曲线光滑”替代。
\end{ideabox}

\section{物理公式基线}

\subsection{场强、对偶场强与规范链}

采用参考文件中的场强约定
\begin{equation}
F_{\mu\nu}=\frac{i}{g}[D_\mu,D_\nu],\qquad
D_\mu=\partial_\mu-igA_\mu,
\qquad
\Tr(T^aT^b)=\frac12\delta^{ab},
\label{eq:field-strength}
\end{equation}
以及格点四叶 Clover 场强
\begin{equation}
F_{\mu\nu}(x)=-\frac{i}{8}\left[Q_{\mu\nu}(x)-Q_{\mu\nu}^{\dagger}(x)\right],
\qquad
F_{\mu\nu}^{\rm tl}=F_{\mu\nu}-\frac{\Tr F_{\mu\nu}}{N_c}\mathbf 1,
\label{eq:clover}
\end{equation}
和 Euclidean 对偶
\begin{equation}
\widetilde F_{\rho\sigma}
=\frac12\epsilon_{\rho\sigma\mu\nu}F_{\mu\nu}.
\label{eq:dual}
\end{equation}
来源：\pathref{pyqcd/renorm/_tmd.py:317-337}、
\pathref{/root/MyQCD/docs/pyqcd_physics/report_gluon_tmd_lattice_qcd_review_20261023.tex:648-690}。

一维 staple 从 \(x\) 到 \(x+z\hat z+b_\perp\) 的代码路径是
\begin{equation}
\Gamma_{z,b}:\quad
-L\hat z\longrightarrow+b_\perp\hat e_\perp
\longrightarrow+(L+z)\hat z,
\qquad L=12a\ \text{（当前生产）}.
\label{eq:staple-path}
\end{equation}
二维 staircase 为
\begin{equation}
\Gamma^{(0,1)}_{2d}:\quad
-L\hat z\to b_x\hat x\to b_y\hat y\to(L+z)\hat z,
\label{eq:staple-2d-01}
\end{equation}
反序 \((1,0)\) 交换两个横向段，但它是另一条连续 contour，不只是同一连续 Wilson line 的格距离散；有限格距的非交换链接进一步放大路径差。

\subsection{核子胶子双线性}

采用颜色闭合
\begin{equation}
M^{\mu\lambda;\nu\rho}(z,b_\perp;\vec x)
=\Tr\!\left[
F^{\nu\rho}(\vec x)\,
W_{\mathrm{st}}(\vec x,\vec x+z\hat z+\bT)\,
F^{\mu\lambda}(\vec x+z\hat z+\bT)\,
W_{\mathrm{st}}^{\dagger}(\vec x,\vec x+z\hat z+\bT)
\right],
\label{eq:tmd-bilocal}
\end{equation}
unpolarized 在两个端点都取 \(F\)；helicity-labelled 通道把\textbf{原点端点}
的第一个场强换成 \(\widetilde F\)，第二个端点仍是 \(F\)：
\begin{equation}
M_H^{\mu\lambda;\nu\rho}
=\Tr\!\left[
\widetilde F^{\nu\rho}(\vec x)\,
W_{\mathrm{st}}(\vec x,\vec x+z\hat z+\bT)\,
F^{\mu\lambda}(\vec x+z\hat z+\bT)\,
W_{\mathrm{st}}^{\dagger}(\vec x,\vec x+z\hat z+\bT)
\right].
\label{eq:helicity-bilocal}
\end{equation}
对偶使用 \(\epsilon_{xyzt}=+1\) 的固定指标顺序；整体取向可翻转 helicity 算符，
但不能改变 \(T_H-S_H\) 的相对号。当前生产在基本表示中计算迹；伴随归一需要
显式因子 2。有限 \(b\) 的 helicity 分支只有内部代数/几何契约，没有独立
有限 \(b\) oracle；\texttt{pol35} 两点 C2 的公开参考一致不能升级为
“有限 \(b\) helicity OPE 已独立验证”。

设 \(t\) 为 Euclidean 时间方向，\(i,j\) 为纵向 \(z\) 之外的另外两个空间方向。
六组件定义为
\begin{align}
ti&=M_{ti;ti},&
tj&=M_{tj;tj},&
ij_{\rm single}&=M_{ij;ij},\\
M_{titi}&=ti+tj,&
M_{ijij}&=2\,ij_{\rm single},&
combined&=M_{titi}\mp M_{ijij}.
\label{eq:six-components}
\end{align}
unpolarized 物理组合使用减号：
\begin{equation}
O_U=M_{titi}-M_{ijij}
=M_{ti;ti}+M_{tj;tj}-2M_{ij;ij}.
\label{eq:unpolarized-combination}
\end{equation}
helicity producer 中历史 \code{combined} 字段保存加号
\(M_{titi}+M_{ijij}\)，而本区间下游物理选择器明确重建
\begin{equation}
O_H^{\rm phys}=M_{titi}-M_{ijij}.
\label{eq:helicity-physical-combination}
\end{equation}
因此不能把历史 \code{combined} 字段直接当作最终物理螺旋度算符。

正、负纵向方向保留为
\begin{equation}
O_{+}=O(+z),\quad O_{-}=O(-z),\quad
O_{\rm even}=O_++O_-,\quad
O_{\rm odd}=O_+-O_-,
\label{eq:orientation}
\end{equation}
不额外乘 \(1/2\)。这保证 \(z=0\) 时 \(O_{\rm odd}=0\) 可被严格检查。

\subsection{source-resolved \texttt{nopol/pol35} C2}

DeGrand-Rossi 基中 \code{gamma(0)} 是单位阵哨兵，\code{gamma(4)} 是时间矩阵。
核子投影采用
\begin{equation}
P_{\rm nopol}=\frac12(\gamma_0+\gamma_4),
\qquad
P_{\rm pol35}=P_{\rm nopol}\,(i\gamma_3\gamma_5).
\label{eq:two-point-projectors}
\end{equation}
因此 \texttt{nopol}/\texttt{pol35} 是同一正宇称子空间中的 projected correlators：
前者无额外自旋投影，后者乘 \(i\gamma_3\gamma_5\)。在非零动量下它们是插值算符
投影约定，不严格隔离所有同动量的宇称/自旋态，不能写成已经完成的物理态投影。
二者都不是“胶子极化”本身；胶子 helicity 由算符中的
\(F\widetilde F\) 表示，核子自旋由 \texttt{pol35} 表示。两种“极化”不能混称。

真实 C2 缓存为
\begin{equation}
C_{2,\Pi}(\ts;t_{\rm src},t_{\rm sink}),\qquad
\Pi\in\{\texttt{nopol},\texttt{pol35}\},
\label{eq:source-resolved-c2}
\end{equation}
轴语义为 \code{(source,dt)}，其中 \code{dt=(t_{\rm sink}-t_{\rm src}) mod N_t}
是 source-relative time，不是 raw sink；当 \(t_{\rm sink}<t_{\rm src}\) 时施加
反周期负号。合同以 complex64 输入计算、以 complex128 累加并存储，
“存储 complex128”不能当作有效精度 complex128。

\subsection{disconnected C3/C2}

对每个通道 \(a\)、源 \(s\)、插入时刻 \(t_{\rm ins}\)、分离 \(\ts\)、
动量 \(P_z\) 和流时 \(\tau\)，组态 \(k\) 上定义
\begin{equation}
C_{3,a}(t_{\rm ins};\ts,P_z,\tau)
=\left\langle\frac1{N_s}\sum_s O_{a,s}(t_{\rm ins})C_{2,a,s}(\ts)\right\rangle_k
-\frac1{N_s}\sum_s
\left\langle O_{a,s}(t_{\rm ins})\right\rangle_k
\left\langle C_{2,a,s}(\ts)\right\rangle_k,
\label{eq:disconnected-ratio}
\end{equation}
\begin{equation}
R_a(t_{\rm ins};\ts,P_z,\tau)
=\frac{C_{3,a}(t_{\rm ins};\ts,P_z,\tau)}{C_{2,\rm den}(\ts,P_z,\tau)}.
\label{eq:ratio-def}
\end{equation}
source 平均在每一次 JK 删除重算内部完成；不同 source 不是独立统计样本。
在当前逐时间片空间求和归一化下 \(R_a\) 可作为无量纲量报告，但 \(C_2\) 中心值
仍是带插值算符和格点归一化的 bare 量，不能跨路径直接比较；谱学关系
\(R\to\langle P|O|P\rangle/(2E_P)\) 尚未用于提取物理矩阵元。
helicity 的 \(C_3\) 协方差使用 \texttt{pol35}，unpolarized 使用
\texttt{nopol}；两个通道当前共享 \texttt{nopol} 分母。物理 helicity 只在
全部统计完成后取
\begin{equation}
A_H=\Real[-iR_H]=\Imag R_H.
\label{eq:helicity-phase}
\end{equation}
不能在真空扣除或 resampling 前把 \texttt{pol35} 截成实部/虚部。

\subsection{delete-one jackknife}

对 \(N\) 个组态，删除第 \(k\) 个组态后重新计算非线性 ratio：
\begin{equation}
R^{(-k)}=\frac{C_3^{(-k)}}{C_{2,\rm den}^{(-k)}}.
\label{eq:jk-ratio}
\end{equation}
JK 偏差修正和实/虚误差为
\begin{align}
R_{\rm bc}&=N R_{\rm central}-(N-1)\overline{R^{(-k)}},\\
\sigma_{\Real/\Imag}
&=\left[
\frac{N-1}{N}\sum_k
\left(\Real/\Imag R^{(-k)}
-\Real/\Imag\overline{R^{(-k)}}\right)^2
\right]^{1/2}.
\label{eq:jk-error}
\end{align}
跨 \(z,b,\ts\) 的协方差必须来自同一组态索引，因此 11 组态的最大正秩是 10。
JK 删除代数已由契约测试验证；11 个规范组态之间的自相关长度和误差是否充分
包含 gauge autocorrelation 没有独立检查，标为未验证。

\subsection{复数分母与实数分母的代数差异}

当前 PyQCD estimator 计算完整复数除法 \(R=C_3/C_2\)；
归档 donghx estimator 使用 \(\Real C_2\) 作分母：
\begin{align}
R_U^{\rm PyQCD}
&=\Real\frac{C_{3,U}}{C_{2,N}},
&
R_H^{\rm PyQCD}
&=\Imag\frac{C_{3,H}}{C_{2,N}},\\
R_U^{\rm donghx}
&=\frac{\Real C_{3,U}}{\Real C_{2,N}},
&
R_H^{\rm donghx}
&=\frac{\Imag C_{3,H}}{\Real C_{2,N}}.
\label{eq:denominator-conventions}
\end{align}
若 \(C_2=D e^{i\phi}\)，则两者通过 \(\tan\phi\) 混合 \(C_3\) 的实部与虚部。
这是一项\textbf{估计量/相位约定差异}，不是紫外算符混合，也不是胶子-夸克 mixing。
相关源码：
\pathref{pyqcd/analysis/_gradient_flow_ratio.py:137-214}、
\pathref{refer/donghx/_Gradient_Flow_GluonPDF/calc_ratio_matched.py:320-372}。

\subsection{梯度流与物理窗口}

连续 Wilson flow 和流半径定义为
\begin{equation}
\partial_t B_\mu(t,x)=D_\nu G_{\nu\mu}(t,x),
\qquad
r_f=\sqrt{8t}=\sqrt{8\tau}\,a.
\label{eq:flow-radius}
\end{equation}
当前使用 \(\tau=3\)，故
\begin{equation}
r_f=4.899a\simeq0.516\,\mathrm{fm}.
\label{eq:current-rf}
\end{equation}
短流时匹配要求
\begin{equation}
a\ll r_f\ll \min(z,b_T,\Lambda_{\rm QCD}^{-1}),
\qquad
\Pz r_f/(\hbar c)\ll1.
\label{eq:flow-window}
\end{equation}
当前所有非零横向距离与 \(z/a=1,\ldots,4\) 均小于 \(r_f\)，且 P200/P400 的 \(\Pz r_f/(\hbar c)\) 均大于 1。
因此最终对象只能是 fixed-flow diagnostic，不能解释为受控短流时 TMD。

\subsection{完整 TMD 链还要求什么}

在 ratio 诊断之后，完整胶子 TMD-PDF 还需要至少以下串联门：
\begin{align}
&\text{UV/端点:}\quad Z_O,\ Z_{ij}^{\rm UV},\
\text{Wilson-line subtraction},\ \sqrt{Z_E}\ \text{(同几何 Wilson loop)},\\
&\text{rapidity/soft:}\quad
f^R=\frac{f^{\rm bare}}{\sqrt{S_R(\bT;\mu,Y_1,Y_2)}},\quad
2\zeta\frac{\partial}{\partial\zeta}\ln f^R=K_R,\\
&\text{matching:}\quad
\widetilde f_i=C_{ij}\otimes f_j,\quad i,j\in\{g,q_S,\bar q_S\},\\
&\qquad C_{ij}=\delta_{ij}+\frac{\alpha_sC_A}{2\pi}g_{ij}+\cdots,\quad
C_{ij}\ \text{含 CS 因子}\ e^{\frac12K_j\ln(\zeta_z/\zeta)},\\
&\text{小 }b_T:\quad \text{OPE},\ b_T^*\ \text{模型},\
O\!\left(\frac{1}{b_T^2P_z^2}\right)\ \text{幂修正},\\
&\text{limits:}\quad L_\parallel\to\infty,\quad
\tau\to0,\quad z,b_T\to0,\quad \Pz\to\infty,\quad a\to0,\quad V\to\infty,
\label{eq:full-tmd-chain}
\end{align}
其中 \(C_{ij}\)（硬匹配）与 \(Z_{ij}^{\rm UV}\)（紫外重正化）不是同一对象；
胶子-夸克单态最低需要 \(g/q_S/\bar q_S\) 基底，不能只保留 gg 子块。
上述极限不可交换（尤其 \(\tau\to0\) 与 \(z,b_T\to0\)）。
当前区间没有完成任何一项的物理闭环；代码接口存在不等于 all11 数据已经消费它们。

\section{十五步交叉工作流：B1--B6}

\subsection{B1 目的与前期准备}

本工作的目的是把“有限横距”从几何参数推进到可重构的统计 observable，并检验
真实核子 C2 与胶子 OPE 是否能在共享组态索引上进入 disconnected ratio。
前置条件是 \(24^3\times72\) 真实蒸馏数据、复现的 perambulator/VVV、规范场流化、
idempotent 投影检查、反周期时间边界和 source-resolved C2 contract。

\subsection{B2 输入}

\begin{xltabular}{\textwidth}{@{}p{0.25\textwidth}p{0.25\textwidth}X@{}}
\caption{B2 主要输入与约束（来源：manifest、cache contract 与源码）}\\
\toprule
输入 & 当前值/形状 & 物理与数值约束 \\
\midrule
\endfirsthead
\toprule
输入 & 当前值/形状 & 物理与数值约束 \\
\midrule
\endhead
组态 & 4150, 6250, 6450, 6650, 6850, 7050, 7250, 7450, 7650, 7850, 8050 & 11 个真实组态；单系综、单格距 \\
格点/流 & \(24^3\times72\)，\(\tau=3\)，\(\epsilon=0.05\) & \(a\simeq0.1053\,\mathrm{fm}\)，\(r_f=4.899a\) \\
OPE 几何 & \(z/a=0,\ldots,12\)，1D \(b/a=0,\ldots,4\) & full channel \(2\times4\times6\)；复数保留 \\
二维几何 & \(b_x/a=0,\ldots,4\)，\(b_y/a=0,1\) & staircase \((0,1)\)，另有 4150/6250 的 \((1,0)\) 对照 \\
核子 C2 & 每动量 \((72,72)\)，c64 计算/c128 存储 & source-resolved，反周期边界，Cg5g4，\texttt{nopol/pol35} \\
动量 & P000/P200/P400 & \(\Pz=0,2,4\) 个单位 \(2\pi/L\) \\
ratio 时间 & \(\ts/a=5,\ldots,10\) & insertion \(0\le i\le\ts\) 有效；不用零填充 \\
\bottomrule
\end{xltabular}

\subsection{B3 输出应该是什么}

单组态 full-channel OPE 应为
\begin{equation}
O:(2,4,6,N_z,N_b,N_t),
\qquad
O_{\rm 2d}:(2,4,6,N_z,N_{bx},N_{by},N_t),
\label{eq:ope-shapes}
\end{equation}
复数、逐时间片空间求和、不提前投影。跨组态 ratio 应为
\begin{align}
R_{1d}&:(2,4,6,1,13,5,6,11,1),\\
R_{2d}&:(2,4,6,1,13,5,2,6,11,1),
\label{eq:ratio-shapes}
\end{align}
分别对应 channel、orientation、component、direction、\(z\)、\(b\) 或
\((b_x,b_y)\)、\(\ts\)、insertion、momentum；JK 样本再增加 11 个 delete-one 组态轴。

\subsection{B4 整体框架}

\begin{table}[htbp]
\centering
\small
\caption{B4 物理链与软件链的对应（数据来源：源码实测）}
\begin{tabularx}{\textwidth}{@{}l X l@{}}
\toprule
物理步骤 & 数学对象 & 软件门 \\
\midrule
流化 & $V_\tau=\mathcal T\exp[\int_0^\tau Z(V_s)\dd s]V_0$ & flow contract \\
场强/路径 & $F_{\mu\nu},W_{\rm st}$ & geometry contract \\
胶子 OPE & $M^{ti;ti},M^{tj;tj},M^{ij;ij}$ & full-channel producer \\
核子 C2 & $P_\Pi C_{2}$ & source-resolved v3 cache \\
真空扣除 & $\langle OC_2\rangle-\langle O\rangle\langle C_2\rangle$ & ratio estimator \\
统计 & delete-one JK & covariance NPZ/JSON \\
物理解释 & plateau + soft/rapidity + matching + limits & 当前 fail-closed \\
\bottomrule
\end{tabularx}
\end{table}

\subsection{B5 关键推导与代码对应}

代码直接实现的核心段有三层。第一层是核子投影和 source-resolved C2：
\lstinputlisting[language=Python,firstline=930,lastline=941]{/root/PyQCD/pyqcd/pipeline/_steps.py}
第二步由同一个 \(Cg5g4\) 矩阵得到两个投影，并保留复数与反周期符号：
\lstinputlisting[language=Python,firstline=1009,lastline=1028]{/root/PyQCD/pyqcd/pipeline/_steps.py}
第三步把六组件和方向组合写入 full-channel OPE：
\lstinputlisting[language=Python,firstline=345,lastline=363]{/root/PyQCD/pyqcd/renorm/_tmd.py}

ratio 的估计量先对 source 求平均、再做组态 covariance；下列代码给出中心 \(C_3\)
和共同分母：
\lstinputlisting[language=Python,firstline=137,lastline=152]{/root/PyQCD/pyqcd/analysis/_gradient_flow_ratio.py}
JK 部分直接使用组态删除后的解析展开：
\lstinputlisting[language=Python,firstline=164,lastline=199]{/root/PyQCD/pyqcd/analysis/_gradient_flow_ratio.py}
\lstinputlisting[language=Python,firstline=200,lastline=209]{/root/PyQCD/pyqcd/analysis/_gradient_flow_ratio.py}

\subsection{B6 正确执行顺序}

\begin{table}[htbp]
\centering
\small
\caption{B6 执行顺序与停止条件（示意，依据源码与 validation reference）}
\begin{tabularx}{\textwidth}{@{}l X@{}}
\toprule
步骤 & 正确做法与失败条件 \\
\midrule
1. 输入 & 固定系综、体积、a、\(\tau\)、\(\epsilon\)、L、方向、动量；身份缺失即失败 \\
2. 流化 & 从 raw gauge 独立流化；监测 \(U^\dagger U\)、det 和流作用量；失败不进入 OPE \\
3. OPE & 同一流场构造 canonical Clover、dual、staple；端点、颜色和 \(b=0\) 极限先通过 \\
4. C2 & 同一合同一次生成 \texttt{nopol/pol35}；保留 complex128、source 轴和反周期符号 \\
5. audit & 检查极化成对、七/八轴 OPE、shape/labels、共享组态和 SHA；不通过则 blocked \\
6. ratio & 仅用同一组态顺序；insertion 越界填 NaN；在每样本重算 C2 和 C3 \\
7. 统计 & 报告 JK rank、条件数、finite fraction、\(\ts\) 漂移；漂移大则 diagnostic \\
8. 物理 & 只有流窗口、平台、soft/rapidity、matching 和极限都通过才允许物理命名 \\
\bottomrule
\end{tabularx}
\end{table}

\section{十五步交叉工作流：B7--B15}

\subsection{B7 实际执行情况}

本区间真实执行由 2026-10-25 的 all11 批次完成。关键实测过程如下：

\begin{xltabular}{\textwidth}{@{}p{0.25\textwidth}X@{}}
\caption{B7 实测过程（来源：all11 final log）}\\
\toprule
阶段 & 命令/动作、退出状态与关键输出 \\
\midrule
\endfirsthead
\toprule
阶段 & 命令/动作、退出状态与关键输出 \\
\endhead
双极化 C2 & 11 组态均生成 P000/P200/P400 的 \texttt{nopol/pol35} \((72,72)\) H5，共 66 个；c64 计算、c128 存储 \\
随机 oracle & 每组态随机时间对的独立合同对照，11 组态最大绝对误差 \(1.63\times10^{-11}\) \\
公共参考 & 4150/P200 对 \texttt{/public} 的 nopol \(\mathrm{rel}_{L2}=3.51\times10^{-8}\)，pol35 \(\mathrm{rel}_{L2}=3.71\times10^{-6}\) \\
1D OPE & 11 个 full-channel OPE，shape \((2,4,6,13,5,72)\)，dtype complex64 \\
2D OPE & 13 个二维 OPE，shape \((2,4,6,13,5,2,72)\)；11 个 order \((0,1)\)，2 个 order \((1,0)\) \\
audit & 1D all11 与 2D all11 均 \code{status=ready}，\code{blocking_reasons=[]}；ready 只表示输入契约 \\
ratio & 1D/2D 三动量、六 \(\ts\)、delete-one JK 均完成；当前状态为 \statusdiag \\
GPU flow & 7650 的 \(72\times24^3\) flow 用时 238 s，峰值 allocated/reserved 3972.1/4462.0 MiB；这是修复后的独立 GPU 验证，不是 all11 全部组态的生产 provenance（部分组态曾 fallback 到 CPU） \\
全量回归 & 本报告再次实测 \code{python -B -m pyqcd.testing}：112 passed, 0 skipped, 0 failed \\
\bottomrule
\end{xltabular}

\subsubsection{B7.1 与归档 donghx 的受控交叉验证}

提交 \code{395dba0} 同时引入 \code{pyqcd/testing/tmd/cross_reference_donghx.py}。
在 \(4^4\times8\)、\(\tau\in\{0,0.1,0.5\}\)、\(\epsilon=0.01\) 的受控输入上，
23 项 required checks 全部通过（\code{overall_status=pass}）：
Wilson flow 的 complex64 相对 \(L_2\le2.5\times10^{-7}\)、直线胶子 OPE
\(\le3.3\times10^{-6}\)、flow-to-quasi 与 gg matching kernel 逐点一致、
同 \(\hbar\) Fourier 差 \(1.39\times10^{-17}\)、C3 中心值差为 0。
默认 \(\hbar\) 常数下的 Fourier 差为 \(1.87\times10^{-4}\)，必须显式固定约定。
唯一实质差异是 denominator：对同一复输入，复分母与 donghx 实分母的
\(\max|\Delta R|=6.77\times10^{-3}\)（归一化 7.86\%，相对 \(L_2\) 7.28\%）。
有限 \(b\) 的 helicity TMD 在归档参考中没有对应实现，因此“23/23”不能
升级为有限 \(b\) 物理链的独立验证。

证据：\pathref{pyqcd/testing/tmd/artifacts/donghx_cross_reference_20261025_082205/report.md:1-31}、
\pathref{pyqcd/testing/tmd/artifacts/donghx_cross_reference_20261025_082205/comparison.json}、
\pathref{logs/v2026102422.txt:29-42}。

\subsection{B8 实际输出情况}

\begin{table}[htbp]
\centering
\small
\caption{B8 结果规模与状态（来源：all11 summary/manifest 实测）}
\begin{tabularx}{\textwidth}{@{}l l l X@{}}
\toprule
对象 & 数值/形状 & finite & 物理解释 \\
\midrule
组态数 & 11 & --- & 仅 JK 最大正秩 10，不构成大样本生产 \\
1D valid ratio & 159120/动量（raw 205920，NaN 46800） & 100\% & 仅 valid-mask 子集有限，不代表独立信息量 \\
1D helicity phase & 3315/动量 & 100\% & 有限性通过，不代表相位平台稳定 \\
2D valid ratio & 318240/动量（raw 411840，NaN 93600） & 100\% & 二维 schema 和轴线通过 \\
2D helicity phase & 6630/动量 & 100\% & 二维相位有限性通过 \\
1D covariance & \(780\times780\) & rank 10 & 最大秩等于 \(N_{\rm conf}-1\) \\
2D covariance & \(1560\times1560\) & rank 10 & 维度远大于统计秩，不能全协方差反演 \\
\bottomrule
\end{tabularx}
\end{table}

来源：
\pathref{data/runs/tmd-production-10conf/analysis/finite_b_ratio_v3_all11/all11_ratio_v3_summary.json:41-151}、
\pathref{data/runs/tmd-production-10conf/analysis/finite_b_ratio_v3_2d_all11/all11_ratio_2d_v3_summary.json:45-165}、
\pathref{data/runs/tmd-production-10conf/analysis/source_resolved_gluon_tmd_all11_manifest_20261025.json:64-108}。

\begin{keybox}
所有 \code{finite_fraction=1} 和 \code{covariance_rank=10} 都只证明“没有 NaN/Inf”以及
“达到样本秩上限”。它们不证明每个 \(z,b,\ts\) 点是独立、稳定或物理可用的。
口径说明：1D 核心函数的 schema 仍是 \code{gradient_flow_gluon_tmd_ratio_v2}；
文件名中的 “ratio-v3/all11” 是批次产物名，真正显式 v3 的是 2D 路径和
\code{all11_ratio_v3} 汇总文件。
\end{keybox}

\subsection{B9 全部图表与物理结果分析}

\subsubsection*{图表边界}

\code{git diff --name-only dev25..HEAD} 的实测结果没有任何
\code{.png/.pdf/.svg} 文件。因此该版本区间原先没有新增成图；已有的
\code{data/runs/tmd-production-10conf/analysis/tmd_ratio/*.png} 9 张图生成于
2026-10-22，早于 \code{dev25}，只能作为前序背景，不能冒充本区间结果。
本报告从当前 NPZ/HDF5/JSON 直接生成 9 张 PDF 图，构成本区间的完整图表集。

\begin{xltabular}{\textwidth}{@{}l p{0.27\textwidth}X@{}}
\caption{本报告新增图表清单（数据来源：all11 NPZ/HDF5/JSON）}\\
\toprule
ID & 文件 & 主要物理内容 \\
\midrule
\endfirsthead
\toprule
ID & 文件 & 主要物理内容 \\
\midrule
\endhead
F1 & \pathref{figures/tmd_gluon_audit_20261025/fig01_tsep_stability_1d.pdf} & 1D \(\ts\) 稳定性，P000/P200/P400 的 \(R_U\) 与 helicity phase \\
F2 & \pathref{figures/tmd_gluon_audit_20261025/fig02_b_dependence_1d.pdf} & 参考 \(\ts=8a\) 下的 \(b\) 依赖与 JK 误差 \\
F3 & \pathref{figures/tmd_gluon_audit_20261025/fig03_z_b_maps_1d.pdf} & 三动量在 \((z,b)\) 平面的诊断热图 \\
F4 & \pathref{figures/tmd_gluon_audit_20261025/fig04_bx_by_slices_2d.pdf} & P200 的二维 \((b_x,b_y)\) staircase 切片 \\
F5 & \pathref{figures/tmd_gluon_audit_20261025/fig05_jk_covariance_spectrum.pdf} & 1D/2D JK 协方差正谱 \\
F6 & \pathref{figures/tmd_gluon_audit_20261025/fig06_plateau_instability.pdf} & \(\ts\) 漂移/误差，平台失败量化 \\
F7 & \pathref{figures/tmd_gluon_audit_20261025/fig07_staircase_order_sensitivity.pdf} & \((0,1)\) 与 \((1,0)\) 路径顺序差 \\
F8 & \pathref{figures/tmd_gluon_audit_20261025/fig08_flow_window_check.pdf} & \(r_f/a=4.899\) 与全部非零 \(b\) 的窗口冲突 \\
F9 & \pathref{figures/tmd_gluon_audit_20261025/fig09_nopol_denominator_phase.pdf} & 复/实 \texttt{nopol} 分母约定敏感性 \\
F10 & \pathref{figures/tmd_gluon_audit_20261025/fig10_denominator_convention.pdf} & 真实 11 组态上 real/complex 分母的差异与 JK 显著性 \\
F11 & \pathref{figures/tmd_gluon_audit_20261025/fig11_flow_window_pilot.pdf} & 11 组态 $\tau=1.0/3.0$ full-channel 流窗口试点 \\
\bottomrule
\end{xltabular}

\subsubsection{B9.1 一维 t\_sep 稳定性}

\begin{figure}[htbp]
\centering
\includegraphics[width=0.92\textwidth,height=0.64\textheight,keepaspectratio]{figures/tmd_gluon_audit_20261025/fig01_tsep_stability_1d.pdf}
\caption{1D 诊断：\(z/a=2,b/a=0\)，对每个 \(\ts\) 的 plateau insertion 窗做平均，误差由同一组态轴上的 delete-one JK 计算。纵轴是无量纲 Euclidean ratio，不是 PDF。}
\end{figure}

\textbf{物理意义.} 若基态已隔离，\(R_U\) 或 helicity phase 应随 \(\ts\) 进入稳定平台；
平台内变化应与 JK 误差同量级。\textbf{对应结果.} P000 的 \(R_U\) 从
\(-0.395\pm0.026\) 漂移到 \(0.137\pm0.164\)；P200 从
\(-0.372\pm0.017\) 到 \(-0.100\pm0.054\)；P400 从
\(-0.763\pm0.084\) 到 \(4.160\pm2.074\)。helicity phase 同样没有平台，
P400 在 \(\ts/a=10\) 达到 \(-3.45\pm4.05\)。\textbf{物理解析.} 这些曲线不是
“有噪声的基态矩阵元”，而是未识别激发态、流污染、分母小或统计不足共同导致的失败信号。

\subsubsection{B9.2 b 依赖与 (z,b) 图}

\begin{figure}[htbp]
\centering
\includegraphics[width=0.90\textwidth,height=0.58\textheight,keepaspectratio]{figures/tmd_gluon_audit_20261025/fig02_b_dependence_1d.pdf}
\caption{参考点 \(\ts/a=8,z/a=2\)。左：unpolarized \(R_U\)；右：helicity phase。误差为 insertion 窗口平均后的 delete-one JK。}
\end{figure}

\begin{figure}[htbp]
\centering
\includegraphics[width=0.98\textwidth,height=0.68\textheight,keepaspectratio]{figures/tmd_gluon_audit_20261025/fig03_z_b_maps_1d.pdf}
\caption{1D \((z,b)\) 诊断图：每列一个动量，上排 unpolarized、下排 helicity phase；固定 \(\ts/a=8\) 并对 insertion 窗 \([2,5]\) 平均。颜色尺度按每幅图最大绝对值居中。}
\end{figure}

\textbf{物理意义.} 在理想短流时窗口内，\(b\) 依赖应由非局域胶子关联决定，并应在
旋转/反射和连续极限下具有受控对称性；固定 \(\ts\) 的图只用于发现异常，不是平台值。
\textbf{对应结果.} P000/P200 在 \(b=0\) 的中心值相近，但 \(b/a=1,2\) 的 JK 误差
已达 \(1.7\)--\(3.4\)，与中心值同量级甚至更大；P400 的误差更大（\(3.4\)--\(5.1\)）。
因此不能按中心值排序，也不能把低 \(b\) 的“量级相近”当作稳定结构。
不同 \(z,b\) 的组合没有形成平滑、可由单一尺度描述的曲面。\textbf{物理解析.}
由于 \(r_f=0.516\,\mathrm{fm}\) 大于所有非零 \(b\le0.421\,\mathrm{fm}\)，
流化会直接平滑待测横向结构；因此在当前固定流时下，\(b\) 依赖不能解释为物理 TMD
分布。该图的价值是证明诊断可以运行并暴露尺度层级失败，而不是提取形状。

\subsubsection{B9.3 二维 staircase 结果}

\begin{figure}[htbp]
\centering
\includegraphics[width=0.90\textwidth,height=0.62\textheight,keepaspectratio]{figures/tmd_gluon_audit_20261025/fig04_bx_by_slices_2d.pdf}
\caption{P200 二维结果：固定 \(\ts/a=8\)、对 insertion 窗 \([2,5]\) 平均；每列对应 \(b_y/a=0,1\)，横轴 \(b_x/a\)，纵轴 \(z/a\)。}
\end{figure}

\textbf{物理意义.} 二维 Wilson line 应保留横向向量的方向和离散路径；在 \(b_y=0\)
时，同一 producer 的二维路径顺序必须精确退化。\textbf{对应结果.} 需要区分三件事：
(i) 同一 2D producer 的 \((0,1)\) 与 \((1,0)\) 在 \(b_y=0\) 时逐元素为零；
(ii) 合成 1D/2D 契约在 \(2\times10^{-12}\) 内通过；
(iii) 真实生产的 1D-v2 与 2D-v3 cache 不是同一实现，11 个组态中 6 个在
\(b_y=0\) 的相对 \(L_2\) 为 \(7.48\times10^{-5}\)--\(7.59\times10^{-5}\)
（最大绝对差约 \(0.035\)--\(0.038\)），5 个逐元素一致。因此不能写“2D 与 1D
输入数值一致”，应写“合成退化契约通过，生产 cache 在 complex64 精度上存在
约 \(7.5\times10^{-5}\) 的差异，尚未证明逐点一致”。\textbf{物理解析.} 二维网格
目前只覆盖 \(b_y/a\le1\)，不足以做旋转/反射连续性或大 \(b\) 尾部分析；
它证明“二维轴可执行”，不能证明连续横向对称性。

\subsubsection{B9.4 JK 协方差与平台指标}

\begin{figure}[htbp]
\centering
\includegraphics[width=0.90\textwidth,height=0.52\textheight,keepaspectratio]{figures/tmd_gluon_audit_20261025/fig05_jk_covariance_spectrum.pdf}
\caption{1D/2D JK 协方差正本征值谱，按最大正本征值归一。每幅谱恰有 10 个正模，对应 \(N_{\rm conf}-1=10\)。}
\end{figure}

\begin{figure}[htbp]
\centering
\includegraphics[width=0.90\textwidth,height=0.52\textheight,keepaspectratio]{figures/tmd_gluon_audit_20261025/fig06_plateau_instability.pdf}
\caption{各动量、各通道的最大 \(\ts\) 漂移与 JK 误差比。红虚线 1 表示“漂移与一个 JK 误差同量级”；所有柱子均远大于 1。}
\end{figure}

\textbf{物理意义.} 协方差秩给出可用于线性组合的独立统计方向数；\(\ts\) 漂移/误差
衡量平台选择是否被统计误差支持。\textbf{指标定义.}
图 6 的柱高是 \(\max_{z,b}\{[\max_t R-\min_t R]/\overline{\sigma_t}\}\)；
它只证明至少一个点失败，不是逐点平台检验。
\textbf{对应结果.} 1D 的正谱秩都是 10（阈值 \(\lambda>10^{-8}\lambda_{\max}\)；
P400 有 21 个原始正本征值 \(>10^{-12}\)，其中 10 个在阈值以上），
正谱条件数为 169.24、115.51、79.52；2D 为 91.45、70.72、50.58。P400 的协方差
最大正本征值达到 \(6.235\times10^3\)（1D）和 \(9.334\times10^3\)（2D），说明少数方向
由强涨落主导。1D unpolarized 漂移比为 7.86、10.17、12.95；2D 为 19.84、8.53、63.57。
helicity phase 的对应最大值也都大于 1。按朴素判据 \(\le1\) 的局部通过率为
1D unpolarized \(22/65\)、\(26/65\)、\(13/65\)，helicity \(43/65\)、\(20/65\)、\(16/65\)；
2D 为 \(43/130\)、\(45/130\)、\(24/130\) 与 \(82/130\)、\(41/130\)、\(28/130\)。
\textbf{物理解析.} 满 JK 秩不等于协方差可逆，因为矩阵维度是 780 或 1560，
而秩只有 10；任何全局相关拟合都必须做 SVD 截断或使用低维投影。
当前没有动量在整个 \((z,b)\) 面上通过均匀平台门；局部通过点仍需相关协方差和
预先声明的选点规则，故整体必须固定为 diagnostic。

\subsubsection{B9.5 staircase 顺序敏感性}

\begin{figure}[htbp]
\centering
\includegraphics[width=0.94\textwidth,height=0.56\textheight,keepaspectratio]{figures/tmd_gluon_audit_20261025/fig07_staircase_order_sensitivity.pdf}
\caption{左：4150/6250 的 full 2D OPE 相对 \(L_2\) 差；右：\(b_y/a=1\) 时 unpolarized combined component 的逐 \((z,b_x)\) 相对差。}
\end{figure}

\textbf{物理意义.} 固定 \(b_x,b_y\) 时，\((0,1)\) 与 \((1,0)\) 是两条不同的
连续 staircase 轮廓，不只是同一连续 Wilson line 的两种格距离散；有限格距的
非交换链接只是让路径依赖更显眼。\textbf{对应结果.} 本报告直接读取两个真实
2D OPE HDF5，定义切片相对 \(L_2\) \(\|A-B\|_2/\|A\|_2\)（以 order \((0,1)\)
的 \(b_y=1\) 切片为分母）：\(b_y=0\) 为零，\(b_y=1\) 对 4150 为
\(1.9543\%\)、对 6250 为 \(1.7935\%\)；完整二维数组为 \(1.156\%\) 与 \(1.076\%\)。
\textbf{物理解析.} 该量级远高于 float32 舍入，是 contour/Wilson-loop 结构；
它是否可由某个高扭度或局域展开描述，需要独立短距 OPE 证明。当前只有一个格距
和两个组态，未验证 \(a\to0\) 行为；在连续外推前不能把反序路径直接平均。
该图支持把 \code{transverse_order} 放入 cache、audit、ratio 身份字段。

\subsubsection{B9.6 流窗口冲突}

\begin{figure}[htbp]
\centering
\includegraphics[width=0.88\textwidth,height=0.46\textheight,keepaspectratio]{figures/tmd_gluon_audit_20261025/fig08_flow_window_check.pdf}
\caption{固定 \(\tau=3\) 的流半径与当前非零 \(b/a\)。\(r_f/a=4.899\)，所有非零 \(b=1,\ldots,4\) 都满足 \(b<r_f\)，即落在 \(r_f>b\) 的失败区内；没有 \(b>r_f\) 的非零点。}
\end{figure}

\textbf{物理意义.} 梯度流尺度必须大于格距以压低 UV，又必须小于待测非局域距离以保留
目标结构；大动量还必须满足 \(\Pz r_f/(\hbar c)\ll1\)。\textbf{对应结果.}
\(r_f=4.899a=0.516\,\mathrm{fm}\)；非零 \(b\) 最大只有
\(0.421\,\mathrm{fm}\)，二维最大模长 \(0.434\,\mathrm{fm}\)；\(z/a=1,\ldots,4\)
同样满足 \(z<r_f\)。令 \(L_s=24a\)、\(L_\parallel=12a\)，则
\(L_\parallel=L_s/2\)，不满足 \(L_\parallel\ll L_s/2\)；且 \(t_s\) 平台失败，
\(t_s\gg1/\Delta E\) 也未验证。P200/P400 的物理动量为约
\(0.981/1.962\,\mathrm{GeV}\)，对应 \(\Pz r_f/(\hbar c)\simeq2.57/5.13\)。
\textbf{物理解析.} 当前数据没有同时满足 \(a\ll r_f\ll b_T,z\)、
\(\Pz r_f/(\hbar c)\ll1\) 与 \(L_\parallel\ll L_s/2\) 的非零 \(b\) 点。
这是物理验收失败，且无法通过改善统计误差修复；必须更换流时、b 范围或系综参数
后重新生产。\textbf{范围限定.} 仓内另有 pre-\code{dev25} 的受控低流时扫描
（\(\tau=0.5,1.0\)，\(r_f/a=2,2.828\)，\(b/a=3,4\)，scalar/real v1 OPE），
部分 \((z,b)\) 可满足流窗口；它们不是本次 full-channel \texttt{nopol/pol35} 链，
不能与 all11 结果混同，也没有闭合本链。

\textbf{\(\tau=1.0\) 全通道试点.} 为检验“更换流时”是否改变结论，本报告在全部
11 个组态上以同一 source-resolved C2 重算 full-channel OPE
（\(\tau=1.0\)、\(\epsilon=0.05\)、staple 12、\(z/a\in\{0,3,4,6,8,10,12\}\)、
\(b/a\in\{0,3,4\}\)；\(r_f/a=2.828\)，故 \(b=3,4\) 与 \(z\ge3\) 均在流半径外）。
在 \(\ts/a=8,z/a=6\)：\(\tau=3\) 的 \(R_U\) 在 \(b=0,3\) 为
\(-0.143\pm0.839\)/\(0.065\pm0.031\)（P200）与
\(-1.641\pm1.723\)/\(0.027\pm0.322\)（P400）；
\(\tau=1\) 的 \(b=3\) 为 \(-0.051\pm0.149\)（P200）与
\(-0.045\pm1.900\)（P400）。两套流时都没有形成随 \(b\) 的单调结构，
窗口外点均与零相容或绝对值很小；图 9 给出同一参考点的 \(b\) 依赖。
更定量地，\(\ts\) 漂移/误差在 \((z,b)\) 上的中位数从 \(\tau=3\) 的
1.68/9.03（P200/P400）降到 \(\tau=1\) 的 1.24/3.44，
但最大值仍为 9.6/21.7 与 9.4/22.8。因此更换流时只能改善典型漂移，
不能消除最坏点，平台失败不是单纯由 \(b<r_f\) 造成；小流时同时带来更大的
UV/离散修正（\(b=0\) 归一化明显移动），不能作为物理解。
该试点 11 个组态、delete-one 秩 10，但仍是单格距、固定 staple、单流时对照，
不能提取物理平台；它把“多流时扫描”从建议变成有数据的路线。
产物：\pathref{data/runs/tmd-flow-window-full-tau1p0/}（OPE cache）与
\pathref{data/runs/tmd-audit-20261025/flow_window_pilot.py}。

\begin{figure}[htbp]
\centering
\includegraphics[width=0.94\textwidth,height=0.44\textheight,keepaspectratio]{figures/tmd_gluon_audit_20261025/fig11_flow_window_pilot.pdf}
\caption{11 组态、\(z/a=6\)、\(\ts/a=8\) 的 \(R_U\) 随 \(b\)：红线为 \(\tau=3\) 的 \(b=0,\ldots,4\)，蓝点为 \(\tau=1\) 的 \(b=0,3,4\)。虚线/点线标出 \(r_f/a=4.899/2.828\)；\(b<r_f\) 的点在流半径内。}
\end{figure}

\subsubsection{B9.7 \texttt{nopol} 分母相位}

\begin{figure}[htbp]
\centering
\includegraphics[width=0.88\textwidth,height=0.50\textheight,keepaspectratio]{figures/tmd_gluon_audit_20261025/fig09_nopol_denominator_phase.pdf}
\caption{11 组态、\(\ts/a=8\) 的 source 平均 \texttt{nopol} C2 复数均值。左图为复平面；右图为 \(|\Imag/\Real|\) 及相位角。}
\end{figure}

\begin{table}[htbp]
\centering
\small
\caption{\texttt{nopol} 分母相位诊断（数据来源：11 组态 v3 C2 HDF5，\(\ts/a=8\)，source 平均）}
\begin{tabular}{@{}l r r r r@{}}
\toprule
动量 & \(\Real C_2\) & \(\Imag C_2\) & 相位角 & \(|\Imag/\Real|\) \\
\midrule
P000 & \(3.608\times10^{-4}\) & \(-3.255\times10^{-6}\) & \(-0.517^\circ\) & 0.902\% \\
P200 & \(4.676\times10^{-5}\) & \(6.013\times10^{-7}\) & \(0.737^\circ\) & 1.286\% \\
P400 & \(1.357\times10^{-7}\) & \(2.117\times10^{-8}\) & \(8.869^\circ\) & 15.604\% \\
\bottomrule
\end{tabular}
\end{table}

\textbf{物理意义.} 在无相位约定歧义的无限统计极限，正宇称 C2 可作为正常化分母；
有限样本下若保留复数分母，除法会把分母相位转入 physical ratio。
\textbf{对应结果.} 表 12 是 \(\ts/a=8\) 的 \textbf{单点} source 平均：
P000/P200 的相位小于 1 度，P400 为 \(8.87^\circ\)（\(|\Imag/\Real|=15.60\%\)）。
跨 \(\ts\) 并不稳定：P400 的 \(|\Imag/\Real|\) 在
\(\ts/a=5,8,10\) 分别约为 2.22\%、15.60\%、121.2\%，
相位约为 \(-1.27^\circ\)、\(8.87^\circ\)、\(50.47^\circ\)（中心值）。
因此 15.6\% 只能标记为单点敏感度，不能当作“P400 代表性相位”。
\textbf{三条 denominator 路径.} (i) ratio-v3 使用复 \(C_2\)；
(ii) 仓内主 \texttt{\_disconnected} 的 \texttt{ensemble\_covariance\_ratio}
使用 \(\Real C_2\)；(iii) 归档 donghx 使用 \(\Real C_2\)。
对同一受控复输入，复/实分母的 \(\max|\Delta R|=6.77\times10^{-3}\)
（归一化 7.86\%，相对 \(L_2\) 7.28\%），这是可以直接比较的敏感度上界。
\textbf{物理解析.} 对 \(C_2=De^{i\phi}\)，一阶展开给出
\begin{equation}
R_U^{\rm PyQCD}\approx R_U^{\rm donghx}
+\frac{\Imag C_{3,U}}{\Real C_{2,N}}\tan\phi,
\qquad
R_H^{\rm PyQCD}\approx R_H^{\rm donghx}
-\frac{\Real C_{3,H}}{\Real C_{2,N}}\tan\phi,
\label{eq:phase-leakage}
\end{equation}
但该式只在 \(|\phi|\ll1\) 时可用；\(\ts/a=10\) 的 \(50^\circ\) 已超出适用范围，
应改用精确复除法。统一前不得判定哪一种 denominator 约定在物理上正确，
也不得把两套中心值直接相减。建议把 denominator convention 变为显式选择，
并在同一 JK 样本上报告 complex-denominator 与 real-denominator 的并列 sensitivity。

\subsubsection{B9.8 复/实分母的真实数据敏感度}

\textbf{实现.} \code{covariance_ratio_estimator} 与
\code{calculate_gradient_flow_tmd_ratios*} 入口新增
\code{denominator_mode="complex"|"real"}：disconnected 分子始终使用
复数 \(C_2\)，real 模式只把归一化分母替换为 \(\Real C_2\)。
用冻结的 11 组态 C2 与 full-channel 1D OPE 重算，complex 模式与生产
\code{*_ratio_v3_diagnostic.npz} 逐元素一致
（P000/P200/P400 的 \(\max|\Delta|=0\)），证明新选项未改变默认结果。

\textbf{解析关系.} 对每个 \((z,b,t_{\rm ins},\ts)\)，
\begin{equation}
R^{\rm real}=R^{\rm complex}\frac{C_2^{\rm nopol}}{\Real C_2^{\rm nopol}}
=R^{\rm complex}\left(1+i\frac{\Imag C_2}{\Real C_2}\right),
\label{eq:denominator-factor}
\end{equation}
所以中心值的相对差严格等于分母相位；真实数据上该恒等式的最大偏差为
\(2.3\times10^{-14}\)。但 JK 误差不是简单倍数：每个 delete-one 样本的
分母相位不同，必须用同一 JK 样本并列重算。

\textbf{结果.} 表 13 与图 10 给出 11 组态、\code{even_sum}/\code{combined}、
insertion 窗平均的 unpolarized 与 helicity 结果。
P000/P200 的约定差在所有 \(\ts\) 都小于 \(0.2\sigma\)，可以忽略；
P400 在 \(\ts/a=8\) 达到约 \(1.0/1.2\sigma\)，在 \(\ts/a=10\) 达到
\(5.7\sigma\)（unpolarized）与 \(3.1\sigma\)（helicity），
相对 \(L_2\) 分别为 105\% 与 43\%。
这是 \(|\Imag/\Real|=121\%\) 的相位在晚 \(\ts\) 放大后的结果，不是新的统计涨落。

\begin{table}[htbp]
\centering
\small
\caption{复/实分母在真实 11 组态上的差异（source-resolved C2 + full-channel OPE；\(\sigma\) 取 complex 模式的 delete-one JK）}
\begin{tabular}{@{}l r r r r r@{}}
\toprule
动量 & \(\ts/a\) & \(|\Imag C_2/\Real C_2|\) (\%) & \(\max|\Delta R|\) & unpol. \(\max|\Delta R|/\sigma\) & helicity \(\max|\Delta R|/\sigma\) \\
\midrule
P000 & 5 & 0.34 & \(2.12\times10^{-4}\) & 0.01 & 0.03 \\
P000 & 8 & 0.90 & \(5.76\times10^{-4}\) & 0.02 & 0.05 \\
P000 & 10 & 2.05 & \(2.77\times10^{-3}\) & 0.08 & 0.09 \\
P200 & 5 & 0.62 & \(2.57\times10^{-4}\) & 0.02 & 0.05 \\
P200 & 8 & 1.29 & \(1.16\times10^{-3}\) & 0.03 & 0.10 \\
P200 & 10 & 1.80 & \(7.25\times10^{-3}\) & 0.10 & 0.18 \\
P400 & 5 & 2.22 & \(6.55\times10^{-3}\) & 0.15 & 0.09 \\
P400 & 8 & 15.60 & \(2.52\times10^{-1}\) & 1.02 & 1.21 \\
P400 & 10 & 121.2 & \(6.78\) & 5.67 & 3.08 \\
\bottomrule
\end{tabular}
\end{table}

\begin{figure}[htbp]
\centering
\includegraphics[width=0.98\textwidth,height=0.46\textheight,keepaspectratio]{figures/tmd_gluon_audit_20261025/fig10_denominator_convention.pdf}
\caption{左：\texttt{nopol} \(C_2\) 的 \(|\Imag/\Real|\) 随 \(\ts\)；中/右：unpolarized 与 helicity 通道中 real 与 complex 分母的 \(\max|\Delta R|/\sigma_{\rm JK}\)。虚线为一个 JK 误差。}
\end{figure}

\textbf{物理解析.} 该结果把“约定选择”从口头警告变成可量化系统：
若统一到 real 分母，P400 晚 \(\ts\) 的中心值会移动超过 5 个 JK 误差，
在同一张图上混合两套约定会产生假的 \(\ts\) 趋势；而 P000/P200 的差异
远小于统计误差，不构成当前主导不确定性。无论选哪一种，\(\ts\) 漂移本身
（B9.1）都远大于该约定差，所以状态仍是 \statusdiag。
下一步必须在生产配置中显式写死 \code{denominator_mode} 并写入
manifest/audit 身份，而不是在分析阶段临时换算。

\subsection{B10 综合分析}

\subsubsection{B10.1 对称性、极限与不变量}

\begin{xltabular}{\textwidth}{@{}p{0.22\textwidth}p{0.20\textwidth}X@{}}
\caption{B10 物理不变量与当前证据（来源：contract tests 与真实数据）}\\
\toprule
检查 & 状态 & 说明 \\
\midrule
\endfirsthead
\toprule
检查 & 状态 & 说明 \\
\midrule
\endhead
规范协变 & 通过 & 2D staple 规范协变残差 \(<2\times10^{-11}\)；纯规范 TMD 约 \(10^{-32}\) \\
\(b=0\) 直线退化 & 通过 & full-channel 与 straight OPE 最大差 \(1.243\times10^{-14}\) \\
\(b_y=0\) 二维退化 & 契约通过 & 同一 2D producer 精确退化；合成 1D/2D 契约通过；真实 1D-v2/2D-v3 有 6/11 组态约 \(7.5\times10^{-5}\) 相对差 \\
\(z=0\) odd 零 & 通过 & helicity odd orientation 精确为零 \\
反周期边界 & 通过 & \(t_{\rm sink}<t_{\rm src}\) 施加负号，C2 oracle 通过 \\
有限 \(b\) helicity & 未验证 & 只有内部代数/几何契约，无独立有限 \(b\) oracle \\
二维路径顺序 & 算符身份 & \(b_y=1\) 切片差约 2\%；连续极限未做 \\
\(\ts\) 平台 & 失败 & 无动量在整个 \((z,b)\) 面通过均匀平台；最坏点 P400 达 \(>10\) \\
流窗口 & 失败 & 非零 \(b<r_f\)、\(z\le4<r_f\)、\(\Pz r_f/(\hbar c)>1\)、\(L_\parallel=L_s/2\) \\
分母约定 & 未统一（已量化） & ratio-v3 复 \(C_2\)；real 模式已显式实现并与生产复算逐元素核对（B9.8）；晚 \(\ts\) 差异达 \(5.7\sigma\) \\
soft/rapidity & 未实现 & 无同几何软函数、CS 核或 rapidity subtraction \\
完整 matching & 未实现 & 只有 gg-only 诊断接口，无 \(g/q_S/\bar q_S\) 单态 \\
连续极限 & 未实现 & 单 \(a\)、单 \(\tau\)、单 \(L_\parallel\)、单系综 \\
\bottomrule
\end{xltabular}

\subsubsection{B10.2 量纲与尺度}

取 \(a=0.1053\,\mathrm{fm}\)，\(\hbar c=0.197327\,\mathrm{GeV\,fm}\)，
空间长度 \(L_s=24a\)。本报告的动量以 \(2\pi/L_s\) 为单位：
\begin{equation}
\Pz(n)=n\frac{2\pi}{La}\hbar c.
\label{eq:pz-physical}
\end{equation}
因此 P200/P400 约为 \(0.981/1.962\,\mathrm{GeV}\)。这高于质子质量，但只有约
\(1.0/2.1\,M_N\)，不构成大动量极限族。更重要的是，固定流半径给出
\(\Pz r_f/(\hbar c)=2.57/5.13\)，说明流化尺度在 boost 后的算符上不可忽略。

体积条件方面，\(L_s=2.527\,\mathrm{fm}\)，\(L_s/2=1.264\,\mathrm{fm}\)；
当前 \(b_{\max}=0.421\,\mathrm{fm}\) 满足 \(b\ll L_s/2\)。但
\(L_\parallel=12a\) 恰好等于 \(L_s/2\)，不满足 \(L_\parallel\ll L_s/2\)，
也未独立扫描；\(t_s\gg1/\Delta E\) 同样未验证。因此有限体积通过不弥补
流窗口失败。

\subsubsection{B10.3 统计秩与信息量}

159120 或 318240 个 finite ratio 是张量元素计数，不是 159120 个独立测量。
它们共享 11 个组态，并在 \(z,b,\ts,\) insertion 上高度相关。任何拟合必须使用
跨 observable 协方差或明确降维。当前 \(780\times780\) 和 \(1560\times1560\)
协方差都只有 10 个正模，直接求逆会把零模放大。正确做法是：
\begin{enumerate}
  \item 先在物理选择器下形成低维 observable；
  \item 对协方差做 SVD，报告截断谱和稳定性；
  \item 用相同 JK 样本传播所有非线性操作；
  \item 将窗口、流时、\(b\) 范围和路径顺序作为系统变化，而不是只做中心值。
\end{enumerate}

\subsubsection{B10.4 当前结果的物理层级}

参考文件采用 L0--L3 分级：L0 为算符/代数，L1 为受控矩阵元与统计，L2 为方案一致性，
L3 为多尺度连续极限。当前可判为：
\begin{itemize}
  \item \textbf{L0：达到。} 场强、staple、颜色闭合、一维/二维退化、cache 和合同均有证据。
  \item \textbf{L1：部分达到。} 真实 source-resolved C2、disconnected C3/C2 与 JK 可运行，
        但基态平台失败，故未完成受控矩阵元提取。
  \item \textbf{L2：未达到。} 无同几何 soft/rapidity、ZR/hybrid 运行链和完整 matching。
  \item \textbf{L3：未达到。} 无多 \(a,\tau,\Pz,L_\parallel\) 和稳定联合外推。
\end{itemize}
因此“软件测试通过”与“物理级别”必须分栏报告，不能把 112 项测试通过升级为 L2/L3。

\subsection{B11 结尾总结}

\begin{conclbox}
本区间最重要的物理成果不是某个数值，而是把有限横距胶子 TMD 从“几何函数”推进到
可追溯的 \(C_2/C_3\)、极化和 JK 协方差诊断，并明确证明当前参数点不能作为物理 TMD。
\(b=0\) 直线极限、同一 producer 的二维退化、规范协变和极化合同已经足够扎实，
可作为下一阶段开发的回归基线；真实 1D/2D cache 的 \(7.5\times10^{-5}\) 差异和
有限 \(b\) helicity 的 oracle 缺项仍需补证。非零 \(b\) 的物理结果则因流窗口和
平台失败而必须暂时冻结。
\end{conclbox}

\subsection{B12 结尾评价}

\textbf{优点.} (i) 每个物理轴都有显式 shape/label/SHA 身份；(ii) C2 合同有公共参考和
逐时间对 oracle；(iii) all11 manifest 明确 fail-closed；(iv) 二维路径顺序没有静默平均；
(v) GPU 内存缺陷通过数值等价 oracle 修复。

\textbf{缺点.} (i) 本次 full-channel all11 链共用单一流时，无法分离流污染
（pre-\code{dev25} 的 scalar 低流时扫描不构成本链）；(ii) 分母 convention
未统一；(iii) 没有稳定 \(\ts\) 平台和低维物理 fit；(iv) 二维 \(b_y\) 范围过小；
(v) 完整 TMD 外层没有运行。下一轮应优先修正物理窗口和估计量一致性，而不是增加图数。

\subsection{B13 补充说明：适用范围与下一步}

\begin{warnbox}
\begin{enumerate}
  \item \textbf{流时扫描：} 对 full-channel 链至少选多个 \(\tau\)，逐点检查
        \(a\ll r_f\ll b_T,z\) 和 \(\Pz r_f/(\hbar c)\ll1\)；当前 \(\tau=3\)
        的单点不足以做流时外推。pre-\code{dev25} 的 scalar 低流时扫描只证明
        部分 \((z,b)\) 可满足窗口，不能替代本链。
  \item \textbf{denominator 决策：} 对 complex 与 real \(\texttt{nopol}\) 分母并列重算，
        在相同 JK 样本上报告 \(R_U,R_H\) 的差和相位灵敏度。
  \item \textbf{平台与态分离：} 增加组态、源数，并用二态/多态或 GEVP 控制激发态；
        在平台稳定前不要发布基态矩阵元。
  \item \textbf{staple 长度与二维方向：} 独立扫描 \(L_\parallel\)，扩大 \(b_y\) 并检验
        旋转/反射；在连续外推前继续把 staircase order 保留为算符身份。
  \item \textbf{完整方案：} 同表示 soft/rapidity、Wilson-line subtraction、ZR/hybrid、
        \(g/q_S/\bar q_S\) matching 和多格距连续极限必须全部进入同一条可复现链。
  \item \textbf{统计独立性：} 11 个规范组态的自相关长度与 JK 误差是否充分，
        当前未验证；config id 间隔不能替代 binning/block-length 检查。
\end{enumerate}
\end{warnbox}

\subsection{B14 参考源}

本报告使用三类来源：(i) 用户指定的 MyQCD 理论参考；(ii) 本仓库当前 HEAD 的代码、
测试、日志、JSON/NPZ/HDF5；(iii) 参考文件中已列出的原始论文。核心论文包括
Lüscher 的 Wilson flow \cite{Luscher2010}、Lüscher--Weisz 微扰分析
\cite{LuscherWeisz2011}、Collins 的 TMD 因子化 \cite{Collins2011}、
胶子准 PDF/LaMET \cite{Ji2013,Ji2014,JiReview2021,ZhangGluonLaMET2019,LiMaQiu2019}、
以及胶子 TMD 软函数/CS 核先例 \cite{ZhangSoft2020,ChuCS2022,ChuSoft2023}。
这些论文提供理论边界，不替代当前数据缺项。

\subsection{B15 附录与附件}

\subsubsection{B15.1 图表闭合清单}

本报告 11 张图全部由真实 all11 结果生成，统一由
\pathref{data/runs/tmd-audit-20261025/make_figures.py} 重生成：
F1--F3 来自 1D plateau 协方差 NPZ；F4 来自 2D plateau 协方差 NPZ
（完整 \(z/a=0,\ldots,12\)）；F5 来自 covariance eigenvalues；
F6 来自 plateau/error；F7 来自 4150/6250 的 order \((0,1)/(1,0)\) 2D OPE HDF5；
F8 是固定 \(\tau=3\) 的解析尺度关系；F9 来自 11 组态 v3 \texttt{nopol} C2 HDF5；
F10 来自 \pathref{data/runs/tmd-audit-20261025/denominator_convention.py} 对同一
C2/OPE 的 real/complex 并列重算（complex 路径与生产 NPZ 逐元素一致）；
F11 来自 \pathref{data/runs/tmd-audit-20261025/flow_window_pilot.py} 的
11 组态 $\tau=1.0/3.0$ 对照（delete-one 秩 10，仅作路线验证）。
任何图中都没有用合成值替代真实数据；只有 F8 是解析尺度检查示意。
图 F5 的秩阈值固定为 \(\lambda>10^{-8}\lambda_{\max}\)，F7 的左右两图分别标注
切片相对 \(L_2\) 与逐 \((z,b_x)\) 相对差，避免与完整二维值混淆。

\subsubsection{B15.2 关键源码直引：二维 staircase}

\lstinputlisting[language=Python,firstline=162,lastline=203]{/root/PyQCD/pyqcd/renorm/_tmd.py}

\subsubsection{B15.3 关键源码直引：最终 helicity 选择}

\lstinputlisting[language=Python,firstline=958,lastline=1001]{/root/PyQCD/pyqcd/analysis/_gradient_flow_ratio.py}

\subsubsection{B15.4 关键源码直引：复数分母校验}

\lstinputlisting[language=Python,firstline=1126,lastline=1131]{/root/PyQCD/pyqcd/analysis/_gradient_flow_ratio.py}

\subsubsection{B15.5 复现命令}

\begin{lstlisting}[language=bash,numbers=none]
python -B -m pyqcd.testing
python -B -m pyqcd.testing._gradient_flow_gluon_contract
python -B -m pyqcd.testing._tmd_physical_contract
python -B -m pyqcd.testing._tmd_su3_geometry_contract
python -B -m pyqcd.testing.tmd.audit_finite_b_inputs \
  --run-dir data/runs/tmd-production-10conf \
  --conf-ids 4150,6250,6450,6650,6850,7050,7250,7450,7650,7850,8050
python -B -m pyqcd.testing.tmd.audit_finite_b_inputs --dimension 2d \
  --run-dir data/runs/tmd-production-10conf \
  --conf-ids 4150,6250,6450,6650,6850,7050,7250,7450,7650,7850,8050
python -B data/runs/tmd-audit-20261025/denominator_convention.py
python data/runs/tmd-audit-20261025/make_figures.py
for i in 1 2 3; do xelatex -interaction=nonstopmode -halt-on-error \
  docs/report_tmd_gluon_audit_20261025.tex; done
\end{lstlisting}
2D audit 的 order 对照使用 \code{--transverse-order 0,1}（4150/6250）。

\subsection{B15.6 证据索引补充}

\begin{xltabular}{\textwidth}{@{}p{0.38\textwidth}X@{}}
\caption{本报告的关键可追溯证据}\\
\toprule
证据 & 内容 \\
\midrule
\endfirsthead
\toprule
证据 & 内容 \\
\midrule
\endhead
\pathref{logs/v2026102508.txt:17-65} & 11 组态真实 C2、GPU、1D/2D ratio、JK、staircase 与 112 项回归 \\
\pathref{data/runs/tmd-production-10conf/analysis/source_resolved_gluon_tmd_all11_manifest_20261025.json:64-108} & 11 组态、1D/2D 维度、\(\ts\) 和 fail-closed 状态 \\
\pathref{data/runs/tmd-production-10conf/analysis/finite_b_ratio_v3_all11/all11_ratio_v3_summary.json:41-151} & 1D shape、finite count、JK shape 和 diagnostic status \\
\pathref{data/runs/tmd-production-10conf/analysis/finite_b_ratio_v3_all11/all11_plateau_covariance.json:32-118} & 1D observable selector、JK rank、条件数和 \(\ts\) 漂移 \\
\pathref{data/runs/tmd-production-10conf/analysis/finite_b_ratio_v3_2d_all11/all11_ratio_2d_v3_summary.json:45-165} & 2D shape、finite count、staircase order 和 diagnostic status \\
\pathref{data/runs/tmd-production-10conf/analysis/finite_b_ratio_v3_2d_all11/all11_plateau_covariance_2d.json:33-86} & 2D JK rank、条件数和平台漂移 \\
\pathref{data/runs/tmd-production-10conf/analysis/gpu_flow_memory_regression_20261025.json:2-35} & OOM 根因、修复实现和真实 7650 峰值显存 \\
\pathref{pyqcd/renorm/_tmd.py:123-203} & 1D/2D staple 路径与二维顺序 \\
\pathref{pyqcd/renorm/_tmd.py:345-457} & 六组件、channel、orientation 和 1D full-channel producer \\
\pathref{pyqcd/pipeline/_steps.py:930-1056} & \texttt{nopol/pol35} 投影、Cg5g4 合同和 source-resolved C2 写缓存 \\
\pathref{pyqcd/analysis/_gradient_flow_ratio.py:93-216} & disconnected C3/C2 与解析 delete-one JK \\
\pathref{pyqcd/analysis/_gradient_flow_ratio.py:456-585} & 显式 1D/2D ratio 轴恢复与 schema \\
\pathref{pyqcd/analysis/_gradient_flow_ratio.py:958-1001} & physical helicity \(M_{titi}-M_{ijij}\) 与最终相位 \\
\pathref{/root/MyQCD/docs/pyqcd_physics/report_gluon_tmd_lattice_qcd_review_20261023.tex:1790-1882} & Wilson flow、SFTX 和流窗口理论参考 \\
\pathref{/root/MyQCD/docs/pyqcd_physics/report_gluon_tmd_lattice_qcd_review_20261023.tex:3777-3862} & 共同尺度窗口、误差矩阵和 L0--L3 验收 \\
\bottomrule
\end{xltabular}

\section{最终结论}

\begin{conclbox}
\textbf{准确结论：}从 \code{dev25} 到当前 HEAD，PyQCD 完成了有限横距胶子 TMD
Euclidean 诊断链的实质性扩展，并在 11 个真实组态上生成了完整的 C2/OPE/ratio/JK
材料；软件与算符合同通过 112 项回归。当前数据只能支持 \statusdiag，
不支持物理 TMD-PDF 或 \(\Delta G(x,\bT)\) 声称。下一阶段的决定性问题不是继续增加
同一 \(\tau,a\) 参数点的图，而是先建立合法流窗口、稳定基态平台和统一的
\texttt{nopol} 分母约定（差异已量化，见 B9.8），再进入 soft/rapidity、
ZR/hybrid、matching 与连续极限。
\end{conclbox}

\begin{thebibliography}{99}

\bibitem{local-log-all11}
PyQCD,
\newblock \textit{all11-polarized-tmd-diagnostic final log},
\newblock \pathref{logs/v2026102508.txt}, 2026-10-25.

\bibitem{Collins2011}
J. Collins,
\newblock \textit{Foundations of Perturbative QCD},
\newblock Cambridge University Press, Cambridge, 2011.

\bibitem{Ji2013}
X. Ji,
\newblock ``Parton physics on a Euclidean lattice,''
\newblock Phys. Rev. Lett. \textbf{110}, 262002 (2013);
arXiv:1305.1539.

\bibitem{Ji2014}
X. Ji,
\newblock ``Parton physics from large-momentum effective field theory,''
\newblock Sci. China Phys. Mech. Astron. \textbf{57}, 1407 (2014);
arXiv:1404.6680.

\bibitem{JiReview2021}
X. Ji, Y. Liu, Y.-S. Liu, J.-H. Zhang, and Y. Zhao,
\newblock ``Large-momentum effective theory,''
\newblock Rev. Mod. Phys. \textbf{93}, 035005 (2021); arXiv:2004.03543.

\bibitem{ZhangGluonLaMET2019}
J.-H. Zhang, X. Ji, A. Schäfer, W. Wang, and S. Zhao,
\newblock ``Accessing gluon parton distributions in large momentum effective theory,''
\newblock Phys. Rev. Lett. \textbf{122}, 142001 (2019); arXiv:1808.10824.

\bibitem{LiMaQiu2019}
Z.-Y. Li, Y.-Q. Ma, and J.-W. Qiu,
\newblock ``Multiplicative renormalizability of quasi-parton operators,''
\newblock Phys. Rev. Lett. \textbf{122}, 062002 (2019); arXiv:1809.01836.

\bibitem{ZhangSoft2020}
Q.-A. Zhang \textit{et al.},
\newblock ``Lattice-QCD calculations of TMD soft function through large-momentum effective theory,''
\newblock Phys. Rev. Lett. \textbf{125}, 192001 (2020); arXiv:2005.14572.

\bibitem{ChuCS2022}
M.-H. Chu \textit{et al.},
\newblock ``Nonperturbative determination of Collins--Soper kernel from quasi transverse-momentum dependent wave functions,''
\newblock Phys. Rev. D \textbf{106}, 034509 (2022); arXiv:2204.00200.

\bibitem{ChuSoft2023}
M.-H. Chu \textit{et al.},
\newblock ``Lattice calculation of the intrinsic soft function and the Collins--Soper kernel,''
\newblock JHEP \textbf{08}, 172 (2023); arXiv:2306.06488.

\bibitem{Luscher2010}
M. Lüscher,
\newblock ``Properties and uses of the Wilson flow in lattice QCD,''
\newblock JHEP \textbf{08}, 071 (2010); arXiv:1006.4518.

\bibitem{LuscherWeisz2011}
M. Lüscher and P. Weisz,
\newblock ``Perturbative analysis of the gradient flow in non-Abelian gauge theories,''
\newblock JHEP \textbf{02}, 051 (2011); arXiv:1101.0963.

\end{thebibliography}

\end{document}
